% Part: second-order-logic
% Chapter: sol-and-sets
% Section: comparing-sets

\documentclass[../../../include/open-logic-section]{subfiles}

\begin{document}

\olfileid[pt]{sol}{set}{cmp}

\olsection{Comparação de conjuntos}

\begin{prop}
A !!{formula} $\lforall[x][(X(x) \lif Y(x))]$ define a relação de inclusão,
O.s.\@, $\Sat{M}{\lforall[x][(X(x) \lif Y(x))]}[s]$ sse
$s(X) \subseteq s(Y)$.
\end{prop}

\begin{prop}
A !!{formula} $\lforall[x][(X(x) \liff Y(x))]$ define a relação de
identidade em conjuntos, O.s.\@,
$\Sat{M}{\lforall[x][(X(x) \liff Y(x))]}[s]$ sse $s(X) = s(Y)$.
\end{prop}

\begin{prop}
A !!{formula} $\lexists[x][X(x)]$ define a propriedade de não ser vazio,
O.s.\@, $\Sat{M}{\lexists[x][X(x)]}[s]$ sse $s(X) \neq \emptyset$.
\end{prop}

Um conjunto~$X$ não é maior que um conjunto~$Y$, $\cardle{X}{Y}$, sse
existe !!a{injective} função $f\colon X \to Y$. Como podemos expressar que
uma função é injetiva e também que seus valores para argumentos em~$X$
pertencem a~$Y$, podemos ainda definir a relação de não ser maior que entre
subconjuntos do domínio.

\begin{prop}
A fórmula
\[
\lexists[u][(\lforall[x][(X(x) \lif Y(u(x)))] \land \lforall[x][\lforall[y][(\eq[u(x)][u(y)] \lif
      \eq[x][y])]])]
\]
define a relação de não ser maior que.
\end{prop}

Dois conjuntos têm o mesmo tamanho, ou são ``equinumerosos'',
$\cardeq{X}{Y}$, sse existe !!a{bijective} função~$f\colon X \to Y$.

\begin{prop}
A fórmula
\begin{multline*}
  \lexists[u][(\lforall[x][(X(x) \lif Y(u(x)))] \land {}\\
    \lforall[x][\lforall[y][(\eq[u(x)][u(y)] \lif \eq[x][y])]]
  \land {} \\\lforall[y][(Y(y) \lif \lexists[x][(X(x)
      \land \eq[y][u(x)])])])]
\end{multline*}
define a relação de ser equinumeroso a.
\end{prop}

Abreviaremos essas !!{formula}s, respectivamente, por $X \subseteq Y$,
$X = Y$, $X \neq \emptyset$, $\cardle{X}{Y}$ e $\cardeq{X}{Y}$. (Isso pode
ser um pouco confuso, pois usamos a mesma notação quando falamos
informalmente dos conjuntos $X$ e $Y$---mas aqui a notação é uma abreviação
de !!{formula}s da lógica de segunda ordem que envolvem as variáveis de
relação unária $X$ e~$Y$.)

\begin{prop}
A !!{sentence} $\lforall[X][\lforall[Y][((\cardle{X}{Y} \land
    \cardle{Y}{X}) \lif \cardeq{X}{Y})]]$ é válida.
\end{prop}

\begin{proof}
A !!{sentence} é satisfeita em !!a{structure}~$\Struct{M}$ se, para
quaisquer subconjuntos $X \subseteq \Domain{M}$ e
$Y \subseteq \Domain{M}$, de $\cardle{X}{Y}$ e $\cardle{Y}{X}$ segue-se
$\cardeq{X}{Y}$. Mas isso vale para \emph{quaisquer} conjuntos $X$ e
$Y$---é o teorema de Schr\"oder--Bernstein.
\end{proof}


\end{document}
